%============================================================
\section{Retirement Model}\label{sec:RetirementModel}
%============================================================
In this section, we build a consumption-saving model with an option to retire. We first
state the problem with a hard maximum over work and retirement, and no device that
smooths that margin. We then summarize the Extreme Value Type I taste shocks of
\citet{Iskhakov2017-jb}, which are the usual smoother. We then replace those shocks
with directed search and an endogenous price of retirement wealth, and derive that
problem in detail.

%------------------------------------------------------------
\subsection{The problem without smoothing}
%------------------------------------------------------------
Agents derive utility from consuming non-durable goods $c_t$ and get disutility from
working. The disutility parameter is denoted by $\delta$. Denoting the choice to retire and
the choice to work by $d_t = 0$ and $d_t = 1$, respectively, and assuming logarithmic utility
from consumption, the utility function is
\begin{equation}
u(c_t, d_t) = \log(c_t) - \delta d_t.
\end{equation}
At the end of each period, households save their end-of-period liquid assets $A_t$ at a
risk-free rate $R$. Working in the current period determines whether the agent receives
income $y_t$ in the following period, so beginning-of-period resources are
$M_t = R A_t + y_t d_{t-1}$. We assume that $y_t = y \eta_t$, where $\eta_t$ follows a mean-one
log-normal distribution, $\eta_t \sim \mathcal{N}(-\sigma_\eta^2/2, \sigma_\eta^2)$.

Retirement is self-absorbing, so the retiree's problem is the ordinary consumption-saving
problem. In Bellman form,
\begin{equation}
V_t^{\Ret}(M_t) = \max_{0 \le c_t \le M_t}
\left\{ \log(c_t) + \beta V_{t+1}^{\Ret}(M_{t+1}) \right\},
\end{equation}
where $M_{t+1} = A_t R$ and $A_t = M_t - c_t$. Using the gothic script $V$ ($\goth{V}$) for the
value of ending the period with assets $A_t$,
\begin{equation}
\goth{V}_t^{\Ret}(A_t) \equiv V_{t+1}^{\Ret}(M_{t+1}(A_t)) = V_{t+1}^{\Ret}(A_t R),
\end{equation}
and the retiree's problem in terms of end-of-period assets is
\begin{equation}
\goth{V}_{t-1}^{\Ret}(A_{t-1}) = \max_{0 \le c_t \le A_{t-1} R}
\left\{ \log(c_t) + \beta \goth{V}_t^{\Ret}(A_t) \right\}.
\end{equation}

A worker who has not yet retired has the option to retire this period. The discrete choice
is a hard maximum of the two conditional values, both evaluated at book resources $M_t$,
\begin{equation}
V_t(M_t) = \max \left\{ V_t^{\Ret}(M_t);\ V_t^{W}(M_t) \right\}.
\end{equation}
The worker who continues solves
\begin{equation}
V_t^{W}(M_t) = \max_{0 \le c_t \le M_t}
\left\{ \log(c_t) - \delta + \beta\, \E_t \big[ V_{t+1}(M_{t+1}) \big] \right\},
\end{equation}
with $A_t = M_t - c_t$ and $M_{t+1} = R A_t + y_{t+1}$. Next period's value is again the hard
maximum, so a kink at the wealth level where the household is indifferent between
retiring and working is inherited by the continuation value. The worker's consumption
function therefore jumps at each such threshold, and those jumps propagate backward.
That sawtooth is the object the rest of this section smooths.

%------------------------------------------------------------
\subsection{Extreme Value Type I taste shocks}
%------------------------------------------------------------
The usual smoother puts i.i.d.\ Extreme Value Type I shocks on the two discrete
alternatives. After observing $M_t$, the household draws $(\zeta_t^{\Ret}, \zeta_t^{W})$ and
chooses
\begin{equation}
V_t(M_t) = \max \Big\{
V_t^{\Ret}(M_t) + \sigma_\zeta \zeta_t^{\Ret},\;
V_t^{W}(M_t) + \sigma_\zeta \zeta_t^{W}
\Big\},
\end{equation}
where $\sigma_\zeta > 0$ is the taste-shock scale. The shocks do not enter the continuous
consumption choice inside either branch. Because the draws are i.i.d.\ Extreme Value
Type I, the expected value has the closed form
\begin{equation}
V_t(M_t) = \sigma_\zeta \log \Big(
\exp\big(V_t^{\Ret}(M_t)/\sigma_\zeta\big)
+ \exp\big(V_t^{W}(M_t)/\sigma_\zeta\big)
\Big),
\end{equation}
and the probability of retiring is the logit weight on $V_t^{\Ret}$. As $\sigma_\zeta \to 0$ this
collapses to the hard maximum above. A larger scale puts weight on the inferior branch
near indifference, so the worker's consumption function is smooth in $M_t$. The results
below show this case with the option arriving for sure and with search and account
dispersion shut down.

%------------------------------------------------------------
\subsection{Search and an endogenous account price}
%------------------------------------------------------------
We now replace the taste shocks with two mechanisms that enter together. Each period
there is a possibility of receiving an option to retire. The likelihood depends on the
search effort $\eff_{t-1}$ exerted in the previous period. Search is costly, with
\begin{equation}
\theta(\eff_t) = \sigma_s \left(\eff_t + (1 - \eff_t) \log (1 - \eff_t)\right),
\end{equation}
so the first units of effort are cheap and further effort is increasingly costly. Utility is
then
\begin{equation}
u(c_t, d_t, \eff_t) = \log(c_t) - \delta d_t - \theta(\eff_t).
\end{equation}
We write $\phi_t$ for the account realization. Contingent on receiving the
option, the household observes $\phi_t$ and decides whether to take the draw---retire,
absorbing, with that wealth---or to deny it and work another period. The realization is
the endogenous price of the nest egg at the moment of claiming. It follows a mean-one
log-normal distribution,
\begin{align}
	\log \phi_t \sim \mathcal{N}\Big(-\frac{\sigma_\phi^2}{2},\,\sigma_\phi^2\Big),
	\qquad \mathbb{E}[\phi_t]=1,
	\label{eq:ret_price_dist}
\end{align}
independent of $\eta_t$ and independent across periods. If the household claims, retirement
begins with wealth
\begin{align}
	\tilde M_t = \phi_t M_t,
	\label{eq:ret_price_Mtilde}
\end{align}
and from that point the agent is an ordinary retiree. If the household denies the draw, it
remains a worker with book resources $M_t$ and the standard budget $A_t = M_t - c_t$; the
realization $\phi_t$ is discarded. The draw is a one-shot, rejectable mark-to-market of the
nest egg---a stochastic liquidation value, annuity conversion factor, or defined-contribution
account realization---and not an ongoing price of consumption in retirement. Once the draw
is accepted, $\phi_t$ is locked into the level of wealth and is not redrawn. A high realization
of $\phi_t$ raises the value of retiring and, conditional on claiming, the resources out of which
the retiree consumes. A low realization has the opposite effect and can be rejected.

The retiree's problem is the one stated above. The value of claiming is that function
evaluated at scaled wealth,
\begin{align}
	V_t^{\Ret}(\tilde M_t)
	= V_t^{\Ret}(\phi_t M_t),
	\label{eq:ret_price_VRet}
\end{align}
with claiming-period consumption $c_t^{\Ret}(\phi_t M_t)$ read off the same policy. The retiree's
decision rule is unchanged. Only the wealth at which it is evaluated depends on $\phi_t$.
There is no separate price-contingent retiree problem, and no need to re-solve the retiree
at each quadrature node of $\phi$: the intensive margin is movement along a single
consumption function.

%------------------------------------------------------------
\subsection{Worker's sub-problem}
%------------------------------------------------------------
Each period, a worker either receives the option to retire or has to continue working.
A worker can increase the likelihood of receiving the option to retire by exerting search effort.
The search effort is made at the end of a period and affects the likelihood of receiving
the option to retire in the following period. When period $t$ starts, an income shock is
realized and beginning-of-period book resources $M_t$ are known. The household then
learns whether the option arrived, which was determined by last period's effort. If it did
not, the agent must work, and $\phi_t$ is irrelevant, because there is no account to claim.
If it did, the agent observes $\phi_t$ and decides whether to claim and retire or to deny the
draw and continue working. The discrete comparison is a hard maximum, contingent on
the realized account price. Having made the discrete choice, the agent makes a consumption
decision. If the agent chooses to retire, they stay as a retiree, because retirement is
self-absorbing. If the agent chooses to continue working, they exert search effort, which
then determines the likelihood of receiving a new draw in the next period.
Fig.~\ref{fig:timing_RetirementModel} outlines the timeline of the realization of the shocks and
the decisions in the worker's problem.

\begin{figure}[h!]
	\centering
	\includegraphics[width=\textwidth]{timeline_RetirementModel.pdf}
	\caption{Timeline of the worker's problem}
	\label{fig:timing_RetirementModel}
\end{figure}

Let the set of possible choice sets for the worker's discrete choice be given by
$\mathcal{D}_t \in \{\{\text{Retire}, \text{Work}\}, \{\text{Work}\}\}$
$\equiv \{\{d_t = 0, d_t = 1\}, \{d_t = 1\}\}$. The first element of
this set is the case when the agent has the option to switch and the second element is
the case when the agent does not have the option to switch and must continue working.
For the case when $\mathcal{D}_t = \{W\}$, i.e.\ the worker does not have the option to switch, the
Bellman equation is given by:
\begin{equation}
V_t^{W}(M_t) = \max_{\substack{0 \le c_t \le M_t \\ 0 \le \eff_t \le 1}}
\left\{ u(c_t, 1, \eff_t) + \beta P_1(\eff_t) \goth{V}_t^1(A_t) + \beta P_2(\eff_t) \goth{V}_t^2(A_t) \right\},
\label{eq:ret_full_VW}
\end{equation}
where $A_t = M_t - c_t$, $\goth{V}_t^1(A_t)$ is the expected value of saving $A_t$ when having the
option to retire in $t + 1$, $P_1(\eff_t) = (1 - \eff_t)p + \eff_t$ is the probability of receiving the
option to retire in $t + 1$, given the search effort $\eff_t$ in the current period, $\goth{V}_t^2(A_t)$
is the expected value of saving $A_t$ when not having the option to retire in $t + 1$, and
$P_2(\eff_t) = (1 - \eff_t)(1 - p)$ is the probability of not receiving the option to retire in
$t+1$, given the search effort $\eff_t$ in the current period. Current $\phi_t$ does not enter
\eqref{eq:ret_full_VW}. A worker who did not claim does not cash the account this period.

For the case when $\mathcal{D}_t = \{R, W\}$, i.e.\ the case when the worker has the option to
switch and has observed the account price, the Bellman equation is given by:
\begin{equation}
V_t(M_t, \{R, W\}, \phi_t) = \max
\left\{ V_t^{\Ret}(\phi_t M_t);\ V_t^{W}(M_t) \right\},
\label{eq:ret_price_discrete}
\end{equation}
or
\begin{equation}
\begin{aligned}
V_t(M_t, \{R, W\}, \phi_t) = \max
\Bigg[
&\max_{0 \le c_t \le \phi_t M_t} \left\{ u(c_t, 0, 0) + \beta \goth{V}_t^{\Ret}(A_t) \right\};\\
&\max_{\substack{0 \le c_t \le M_t \\ 0 \le \eff_t \le 1}}
\left\{ u(c_t, 1, \eff_t) + \beta P_1(\eff_t) \goth{V}_t^1(A_t) + \beta P_2(\eff_t) \goth{V}_t^2(A_t) \right\}
\Bigg].
\end{aligned}
\label{eq:ret_full_discrete}
\end{equation}
A high $\phi_t$ makes claiming more attractive and, if accepted, delivers higher retirement
consumption along $c_t^{\Ret}(\cdot)$. A low $\phi_t$ makes claiming less attractive and is denied
whenever $V_t^{\Ret}(\phi_t M_t) < V_t^{W}(M_t)$.

The expected value of saving $A_t$ when having the option to retire in $t + 1$, i.e.\ $\goth{V}_t^1(A_t)$,
now integrates over next period's account draw as well as income,
\begin{equation}
\goth{V}_t^1(A_t) = \int_{y_{t+1}} \int_{\phi_{t+1}}
V_{t+1}\bigl(M_{t+1}(A_t), \{R, W\}, \phi_{t+1}\bigr)\, dF(y_{t+1})\, dF(\phi_{t+1}),
\label{eq:ret_full_gothV1}
\end{equation}
where $M_{t+1} = A_t R + y_{t+1}$, because the agent worked this period, and
$V_{t+1}(M_{t+1}, \{R, W\}, \phi_{t+1})$ is the hard maximum \eqref{eq:ret_full_discrete} one
period ahead. Similarly, the expected value of saving $A_t$ when not having the option to
retire in $t + 1$, i.e.\ $\goth{V}_t^2(A_t)$, is given by
\begin{equation}
\goth{V}_t^2(A_t) = \int_{y_{t+1}}
V_{t+1}^{W}\bigl(M_{t+1}(A_t)\bigr)\, dF(y_{t+1}).
\label{eq:ret_full_gothV2}
\end{equation}
This continuation does not integrate over $\phi_{t+1}$, because a worker who must work
cannot claim. Future draws still sit inside $V^{W}$ through $\goth{V}^{1}$ and $\goth{V}^{2}$ one
period ahead. Equivalently, the beginning-of-period value of a worker who will have the
option, after integrating out the account price, is
\begin{align}
	\bar V_t^{R,W}(M_t)
	\equiv \mathbb{E}_{\phi_t}\Big[
		\max\big\{
			V_t^{\Ret}(\phi_t M_t),\;
			V_t^{W}(M_t)
		\big\}
	\Big],
	\label{eq:ret_full_barVRW}
\end{align}
so that $\goth{V}_{t-1}^{1}(A_{t-1}) = \mathbb{E}_{\eta_t}\bigl[\bar V_t^{R,W}(R A_{t-1} + y_t)\bigr]$. Integrating
\eqref{eq:ret_price_discrete} over $\phi_t$ replaces a single wealth threshold with a family of
thresholds, one for each realization, so the ex-ante claiming probability is interior.

%------------------------------------------------------------
\subsection{Optimal effort and consumption functions}
%------------------------------------------------------------
Since we assume that the consumption decision is made before the effort decision,
\eqref{eq:ret_full_VW} can be split into two sequential sub-period problems and rewritten as:
\begin{equation}
V_t^{W}(M_t) = \max_{0 \le c_t \le M_t} \left[ \goth{V}_t(A_t) \right]
\end{equation}
where
\begin{equation}
\goth{V}_t(A_t) = \max_{0 \le \eff_t \le 1}
\left\{ u(c_t, 1, \eff_t) + \beta P_1(\eff_t) \goth{V}_t^1(A_t) + \beta P_2(\eff_t) \goth{V}_t^2(A_t) \right\}.
\end{equation}
Search effort is chosen only on the working branch and does not depend on the current
draw. The current $\phi_t$ is already sunk from the point of view of next period's option.
Thus, the first order condition gives the optimal effort as a function of the end-of-period
assets that are determined after the optimal consumption decision has been made in the
previous sub-period. The optimal effort is given by
\begin{align}
\frac{\partial \theta(\eff_t)}{\partial \eff_t}
&= \beta \goth{V}_t^1(A_t) \frac{\partial P_1(\eff_t)}{\partial \eff_t}
+ \beta \goth{V}_t^2(A_t) \frac{\partial P_2(\eff_t)}{\partial \eff_t} \\
-\sigma_s \log(1 - \eff_t^*)
&= \beta (1 - p) \goth{V}_t^1(A_t) - \beta (1 - p) \goth{V}_t^2(A_t) \\
\eff_t^* &= 1 - \exp\left( -\frac{\beta (1 - p)}{\sigma_s}
\left( \goth{V}_t^1(A_t) - \goth{V}_t^2(A_t) \right) \right).
\label{eq:ret_full_effort}
\end{align}
This expression states that optimal effort is an increasing function of the excess value
of receiving the option to switch in the next period over not receiving the option to
switch. Account-value dispersion raises $\goth{V}_t^1$ relative to $\goth{V}_t^2$, because the
option to retire is also an option on $\phi_{t+1}$: the household claims when the draw is
attractive and otherwise continues working. That option value increases the excess value
of receiving the option, and therefore increases search effort. The two frictions are
complements. Search governs whether a draw arrives; $\phi_t$ is the content of the draw,
and the hard maximum \eqref{eq:ret_full_discrete} is the decision to take it or to work
another period. Together they smooth both the arrival of the option and the terms on
which it is exercised.

As $\sigma_\phi \to 0$, $\phi_t \to 1$ almost surely and
\eqref{eq:ret_full_VW} is the search problem with a hard maximum at book resources.
Shutting search down as well recovers the single wealth threshold of the unsmoothed
problem.

Next, we derive the expressions for the optimal consumption functions for the retiree
and the worker. Since retirement is self-absorbing, and $\phi_t$ is locked in at the moment
of claiming, the retiree's consumption function is the standard one, evaluated at claimed
wealth:
\begin{equation}
u'(c_t, 0, 0) = R \beta u'(c_{t+1}, 0, 0)
\end{equation}
or
\begin{equation}
c_t^{\Ret}(\phi_t M_t) = \left( R \beta \left( c_{t+1}^{\Ret}(A_t R) \right)^{-1} \right)^{-1},
\end{equation}
where $A_t = \phi_t M_t - c_t^{\Ret}(\phi_t M_t)$. To derive the consumption function for the
worker, we start with the solution to the problem of a worker who does not have the
option to retire, i.e.\ $\mathcal{D}_t = \{W\}$. The first-order condition for
\eqref{eq:ret_full_VW}, with respect to $c_t$, is given by
\begin{equation}
u'(c_t, 1, \eff_t^*) = \beta \left[
P_1(\eff_t^*) \frac{\partial \goth{V}_t^1(A_t)}{\partial A_t}
+ P_2(\eff_t^*) \frac{\partial \goth{V}_t^2(A_t)}{\partial A_t}
\right],
\end{equation}
where $\eff_t^*$ is the optimal effort function in \eqref{eq:ret_full_effort}, and the derivatives
of the end-of-period values are given by
\begin{equation}
\frac{\partial \goth{V}_t^1(A_t)}{\partial A_t}
= R\, \E_t \left\{ \frac{\partial V_{t+1}(M_{t+1}, \{R, W\}, \phi_{t+1})}{\partial M_{t+1}} \right\}
\end{equation}
and
\begin{equation}
\frac{\partial \goth{V}_t^2(A_t)}{\partial A_t}
= R\, \E_t \left\{ \frac{\partial V_{t+1}^{W}(M_{t+1})}{\partial M_{t+1}} \right\},
\end{equation}
where the expectation in the first line is over next period's income and account price, the
expectation in the second line is over next period's income only, and $M_{t+1} = A_t R + y_{t+1}$.
Substituting these expressions yields
\begin{equation}
u'(c_t, 1, \eff_t^*) = R \beta \left[
P_1(\eff_t^*) \E_t \left\{ \frac{\partial V_{t+1}(M_{t+1}, \{R, W\}, \phi_{t+1})}{\partial M_{t+1}} \right\}
+ P_2(\eff_t^*) \E_t \left\{ \frac{\partial V_{t+1}^{W}(M_{t+1})}{\partial M_{t+1}} \right\}
\right].
\end{equation}
The envelope condition on the working branch is
\begin{equation}
\frac{\partial V_t^{W}(M_t)}{\partial M_t} = u'(c_t, 1, \eff_t^*),
\end{equation}
which can be iterated forward to yield
\begin{equation}
\frac{\partial V_{t+1}^{W}(M_{t+1})}{\partial M_{t+1}} = u'(c_{t+1}, 1, \eff_{t+1}^*).
\end{equation}
For the worker who has the option to retire, the value function is the maximum of the
value of claiming and the value of continuing to work,
\begin{equation}
V_t(M_t, \{R, W\}, \phi_t) = \max \left\{ V_t^{\Ret}(\phi_t M_t);\ V_t^{W}(M_t) \right\}.
\end{equation}
An extra unit of book wealth is worth $\phi_t$ units of retirement wealth if the household
claims. Therefore,
\begin{equation}
\frac{\partial V_{t+1}(M_{t+1}, \{R, W\}, \phi_{t+1})}{\partial M_{t+1}}
= \mathbbm{1}_{t+1}^0 \cdot \phi_{t+1} \cdot u'(c_{t+1}, 0, 0)
+ \mathbbm{1}_{t+1}^1 \cdot u'(c_{t+1}, 1, \eff_{t+1}^*),
\end{equation}
where the retiree marginal utility is evaluated at claimed wealth $\phi_{t+1} M_{t+1}$,
$\mathbbm{1}_{t+1}^0 \equiv \mathbbm{1}_{t+1}^0(M_{t+1}, \phi_{t+1}) = 1$ if
$V_{t+1}^{\Ret}(\phi_{t+1} M_{t+1}) \ge V_{t+1}^{W}(M_{t+1})$ and $0$ otherwise, and
$\mathbbm{1}_{t+1}^1 = 1 - \mathbbm{1}_{t+1}^0$. Substituting the two envelopes into the
first-order condition yields the consumption function for the worker who continues to work
\begin{equation}
\begin{aligned}
u'(c_t, 1, \eff_t^*) = R \beta \Big[
&P_1(\eff_t^*) \int_{y_{t+1}} \int_{\phi_{t+1}}
\left[ \mathbbm{1}_{t+1}^0 \cdot \phi_{t+1} \cdot u'\bigl(c_{t+1}^{\Ret}(\phi_{t+1} M_{t+1}), 0, 0\bigr)
+ \mathbbm{1}_{t+1}^1 \cdot u'\bigl(c_{t+1}^{W}(M_{t+1}), 1, \eff_{t+1}^*\bigr) \right]
dF(y_{t+1})\, dF(\phi_{t+1}) \\
&+ P_2(\eff_t^*) \int_{y_{t+1}}
u'\bigl(c_{t+1}^{W}(M_{t+1}), 1, \eff_{t+1}^*\bigr)
\, dF(y_{t+1})
\Big],
\end{aligned}
\end{equation}
or, with logarithmic utility,
\begin{equation}
\begin{aligned}
c_t^{W} = \Bigg( R \beta \Big[
&P_1(\eff_t^*) \int_{y_{t+1}} \int_{\phi_{t+1}}
\left[ \mathbbm{1}_{t+1}^0 \, \phi_{t+1} \left( c_{t+1}^{\Ret}(\phi_{t+1} M_{t+1}) \right)^{-1}
+ \mathbbm{1}_{t+1}^1 \left( c_{t+1}^{W}(M_{t+1}) \right)^{-1} \right]
dF(y_{t+1})\, dF(\phi_{t+1}) \\
&+ P_2(\eff_t^*) \int_{y_{t+1}}
\left( c_{t+1}^{W}(M_{t+1}) \right)^{-1}
\, dF(y_{t+1})
\Big] \Bigg)^{-1},
\end{aligned}
\label{eq:ret_worker_c}
\end{equation}
where $M_{t+1} = A_t R + y_{t+1}$. Because the household observes $\phi_t$ before choosing,
the realized policy of a worker who has the option is the hard choice at that draw. The
ex-ante consumption function is the $\phi$-weighted average of those choices. Working-branch
consumption $c_t^{W}$ and effort $\eff_t^*$ do not depend on the current draw. The envelope
conditions and the EGM algorithm for this baseline are derived in
Appendix~\ref{subsec:appdx_ret_full}.

%------------------------------------------------------------
\subsection{Results}
%------------------------------------------------------------
Having established the solution of the retirement problem, we show how the consumption
rule of the worker who continues to work responds to Extreme Value Type I taste shocks,
to search, and to dispersion in the retirement account. Every figure is drawn at period
$t = 15$ of a $T = 20$ problem, with $R = 1$, $\beta = 0.98$, and $y = 20$, and with no
income uncertainty. The object plotted is the worker's consumption function $C(m_t)$.

With the option arriving for sure ($p = 1$), and with search and account dispersion shut
down, Fig.~\ref{fig:ev} varies the Extreme Value scale $\sigma_\zeta$. At $\sigma_\zeta = 0$
the hard maximum leaves a sequence of sharp discontinuities. Raising the scale to
$0.05$ rounds the peaks. At $0.10$ and $0.15$ the kinks are essentially gone. Larger
taste shocks also shift consumption down through the middle of the resource range and
up at high resources.

%---- EV ----
\begin{figure}[htbp]
\centering
    \includegraphics[width=0.72\linewidth]{EV_VaryT.pdf}
\caption{Worker consumption at $t = 15$ under Extreme Value Type I taste shocks.
The option arrives with certainty ($p = 1$). Search and account dispersion are shut
down ($\sigma_s = 0$, $\sigma_\phi = 0$).}
\label{fig:ev}
\end{figure}

Fig.~\ref{fig:search} shuts the taste shocks and the account price down and sets
$p = 0$, so a worker receives the option only by searching. At $\sigma_s = 0$ the
discontinuities remain. A small search cost, $\sigma_s = 0.15$, trims them from the
right. At $\sigma_s = 0.50$ and $\sigma_s = 1$ the rule is much smoother, and the
smoothing sits to the right of each kink: workers with low liquid resources do not
search, so the left side of each discontinuity is left intact. Higher search costs
also shift consumption up at high resources.

%---- Search ----
\begin{figure}[htbp]
\centering
    \includegraphics[width=0.72\linewidth]{Search_VaryS.pdf}
\caption{Worker consumption at $t = 15$ under search alone. Taste shocks and account
dispersion are shut down ($\sigma_\zeta = 0$, $\sigma_\phi = 0$), and $p = 0$.}
\label{fig:search}
\end{figure}

Account dispersion on its own is shown in Fig.~\ref{fig:price}. Search is shut down
and $p = 1$, so a draw of $\phi_t$ arrives with certainty. In panel~(a), with
$N_\phi = 151$ quadrature nodes, $\sigma_\phi = 0$ reproduces the sawtooth. A small
dispersion, $\sigma_\phi = 0.05$, removes every kink but the first. At $\sigma_\phi = 0.10$
and $0.15$ the rule is smooth, and it lies below the hard-maximum rule through the
middle of the grid. Panel~(b) holds $\sigma_\phi = 0.1$ and varies the number of
nodes. A single node is the hard maximum. Nine nodes follow the smooth rule but leave
a sawtooth of their own. At $N_\phi = 151$ and $N_\phi = 351$ the two curves lie on
top of each other, so further nodes do not change the consumption function.

%---- Price ----
\begin{figure}[htbp]
\centering
\begin{subfigure}{0.48\textwidth}
    \includegraphics[width=\linewidth]{Price_VaryPhi.pdf}
    \caption{Varying account dispersion, $N_\phi = 151$.}
\end{subfigure}
\hfill
\begin{subfigure}{0.48\textwidth}
    \includegraphics[width=\linewidth]{Price_VaryPhiCount.pdf}
    \caption{Varying the number of account-price nodes, $\sigma_\phi = 0.1$.}
\end{subfigure}
\caption{Worker consumption at $t = 15$ with account dispersion alone. Search and
taste shocks are shut down ($\sigma_s = 0$, $\sigma_\zeta = 0$) and $p = 1$.}
\label{fig:price}
\end{figure}

Search and the account price together are shown in Fig.~\ref{fig:searchphi}, again
with $p = 0$ and no taste shocks. In panel~(a), $\sigma_\phi = 0.1$ and
$N_\phi = 151$ are held fixed while $\sigma_s$ rises from $0$ to $1$. The price draw
has already removed most of the sawtooth; raising the search cost removes the
remaining bump at low resources and shifts the rule up. In panel~(b), $\sigma_s = 0.5$
and $N_\phi = 151$ are held fixed while $\sigma_\phi$ rises. At $\sigma_\phi = 0$ the
one-sided kinks of search alone are still visible. They shrink at $\sigma_\phi = 0.05$
and are gone at $0.10$ and $0.15$. Used together, the two mechanisms produce a smooth
consumption function at scales at which either one, used alone, still leaves kinks.

%---- Search and price ----
\begin{figure}[htbp]
\centering
\begin{subfigure}{0.48\textwidth}
    \includegraphics[width=\linewidth]{SearchPhi_VaryS.pdf}
    \caption{Varying search, $\sigma_\phi = 0.1$, $N_\phi = 151$.}
\end{subfigure}
\hfill
\begin{subfigure}{0.48\textwidth}
    \includegraphics[width=\linewidth]{SearchPhi_VaryPhi.pdf}
    \caption{Varying account dispersion, $\sigma_s = 0.5$, $N_\phi = 151$.}
\end{subfigure}
\caption{Worker consumption at $t = 15$ with search and account dispersion together.
Taste shocks are shut down ($\sigma_\zeta = 0$) and $p = 0$.}
\label{fig:searchphi}
\end{figure}

%============================================================
\subsection{Algorithm}
%------------------------------------------------------------
This section outlines the algorithm for solving the problem of the worker who continues
to work in the current period. The retiree is solved first, by the standard consumption-saving
EGM, and is evaluated later at claimed wealth $\phi M$. Taste shocks are shut down, so the
option comparison is a hard maximum. Values are stored in the inverse form
$V^{T} = -1/V$ and recovered by $V = -1/V^{T}$ before that comparison.
After the EGM step, the upper envelope of \citet{Iskhakov2017-jb} drops any decreasing
segments of the endogenous grid.

\begin{algorithm}[htbp]
\caption{EGM-step with search and account dispersion for the worker who continues to work}
\DontPrintSemicolon
Let $\vec{\eta} = \{\eta_1, \dots, \eta_J\}$ and $\vec{\phi} = \{\phi_1, \dots, \phi_K\}$ be quadrature nodes for income and the account price, with weights
$\vec{\omega} = \{\omega_1, \dots, \omega_J\}$ and $\vec{\mu} = \{\mu_1, \dots, \mu_K\}$.\;
Form an ascending grid over end-of-period wealth, $\vec{A}_t = \{A_t^1, \dots, A_t^I\}$.\;
\For{$i = 1, \dots, I$}{
    \For{$j = 1, \dots, J$}{
        Compute book resources $M_{t+1}^j(A_i) = R A_i + y \eta_j$.\;
        Interpolate the worker at book resources:
        $c_{t+1}^{W}(M_{t+1}^j)$ and $V_{t+1}^{W} = -1/V_{t+1}^{T,W}(M_{t+1}^j)$.
        The worker does not depend on $\phi$.\;
        \For{$k = 1, \dots, K$}{
            Compute claimed wealth $\tilde M = \phi_k M_{t+1}^j(A_i)$.\;
            Interpolate the retiree at claimed wealth:
            $c_{t+1}^{\Ret}(\tilde M)$ and $V_{t+1}^{\Ret} = -1/V_{t+1}^{T,\Ret}(\tilde M)$.\;
            Set $V_{t+1}^{RW}(M_{t+1}^j, \phi_k) = \max\{V_{t+1}^{\Ret}(\tilde M),\, V_{t+1}^{W}(M_{t+1}^j)\}$.\;
            Set $\mathbbm{1}_{t+1}^0 = 1$ if $V_{t+1}^{\Ret}(\tilde M) \ge V_{t+1}^{W}(M_{t+1}^j)$, and $0$ otherwise;
            set $\mathbbm{1}_{t+1}^1 = 1 - \mathbbm{1}_{t+1}^0$.\;
            Set $\mathrm{RHS}_1^{\mathrm{pre}}(M_{t+1}^j, \phi_k) = \mathbbm{1}_{t+1}^0 \cdot \phi_k \cdot u'(c_{t+1}^{\Ret}(\tilde M)) + \mathbbm{1}_{t+1}^1 \cdot u'(c_{t+1}^{W}(M_{t+1}^j))$.\;
        }
        Average over the account price:
        $V_{t+1}^{RW}(M_{t+1}^j) = \sum_k \mu_k V_{t+1}^{RW}(M_{t+1}^j, \phi_k)$ and
        $\mathrm{RHS}_1^{\mathrm{pre}}(M_{t+1}^j) = \sum_k \mu_k \mathrm{RHS}_1^{\mathrm{pre}}(M_{t+1}^j, \phi_k)$.\;
        Set $\mathrm{RHS}_2^{\mathrm{pre}}(M_{t+1}^j) = u'(c_{t+1}^{W}(M_{t+1}^j))$. The must-work branch is not averaged over $\phi$.\;
    }
    Compute $\goth{V}_t^1(A_i) = \sum_{j=1}^{J} \omega_j \, V_{t+1}^{RW}(M_{t+1}^j)$ and
    $\goth{V}_t^2(A_i) = \sum_{j=1}^{J} \omega_j \, V_{t+1}^{W}(M_{t+1}^j)$.\;
    Set effort from the two gothic values.
    If $p = 1$, set $\eff_t(A_i) = 0$.
    If $p < 1$ and $\sigma_s = 0$, set $\eff_t(A_i) = 1$.
    Otherwise
    $\eff_t(A_i) = 1 - \exp\bigl(-\beta(1-p)(\goth{V}_t^1(A_i) - \goth{V}_t^2(A_i))/\sigma_s\bigr)$,
    and replace a negative value by $0$.\;
    Compute $\theta(\eff_t(A_i)) = \sigma_s\bigl(\eff_t(A_i) + (1-\eff_t(A_i))\log(1-\eff_t(A_i))\bigr)$,
    with $\theta = \sigma_s$ if $\eff_t = 1$.\;
    Compute $P_1(\eff_t) = (1-\eff_t)\,p + \eff_t$ and $P_2(\eff_t) = (1-\eff_t)\,(1-p)$.\;
    Compute $\mathrm{RHS}_1(A_i) = P_1 \sum_j \omega_j \mathrm{RHS}_1^{\mathrm{pre}}(M_{t+1}^j)$ and
    $\mathrm{RHS}_2(A_i) = P_2 \sum_j \omega_j \mathrm{RHS}_2^{\mathrm{pre}}(M_{t+1}^j)$.\;
    Compute $\mathrm{RHS}(A_i) = R\bigl(\mathrm{RHS}_1(A_i) + \mathrm{RHS}_2(A_i)\bigr)$.\;
    Compute $\E V_{t+1}(A_i) = P_1\,\goth{V}_t^1(A_i) + P_2\,\goth{V}_t^2(A_i)$.\;
    Compute $c_t^{W}(A_i) = u'^{-1}\bigl(\beta\,\mathrm{RHS}(A_i)\bigr)$ and
    $M_t(A_i) = A_i + c_t^{W}(A_i)$.\;
    Compute $V_t^{W}(M_t(A_i)) = u(c_t^{W}(A_i), 1, \eff_t(A_i)) + \beta\,\E V_{t+1}(A_i)$, and store $V_t^{T,W} = -1/V_t^{W}$.\;
}
On the endogenous grid, fill the liquidity-constrained segment below the smallest $M_t(A_i)$ by setting consumption equal to resources.\;
Apply the upper envelope to $(M_t, c_t^{W}, V_t^{T,W})$ and interpolate the consumption and inverse-value functions on the common $M$-grid.\;
\end{algorithm}